\metatitle{Dictionary}
\metadescription{A short glossary of words and terminology that appear in
algebraic combinatorics and symmetric functions.}


\section[wordList]{Dictionary}

This page is a quick glossary for terms that appear across the site.  Most
entries point to a fuller page where the term is used in context.


\subsection[dictionaryAlgebraRepresentation]{Algebra and representation theory}

\textbf{Character.}
The \hyperref[representationCharacterDefinition]{character} of a
representation is the trace of each group element acting on the representation
space.  For $\symS_n$, the
\hyperref[frobeniusCharacteristic]{Frobenius characteristic} converts
characters into symmetric functions.

\textbf{Cluster algebra.}
A cluster algebra is a commutative algebra generated from overlapping sets of
generators called clusters, related by mutation rules.  Cluster structures
often appear in \hyperref[totalPositivity]{total positivity}, Grassmannians,
and canonical-basis questions.

\textbf{Graded character.}
A graded character records the character of a graded representation degree by
degree.  In symmetric-function notation, this often gives a series such as
\(\sum_d q^d\frobChar(M_d)\).

\textbf{Hilbert series.}
The \hyperref[hilbertSeries]{Hilbert series} of a graded vector space records
the dimensions of its homogeneous components.  If the graded space carries an
$\symS_n$-action, the Frobenius characteristic refines the Hilbert series.

\textbf{Quiver.}
A quiver is a directed graph.  Quivers enter this site mostly through
\hyperref[algebras]{quiver Hecke algebras}, Hall--Littlewood-type formulas,
and representation-theoretic models.


\subsection[dictionaryCommutingVariables]{Commuting and anticommuting variables}

\textbf{Bosonic variables.}
Bosonic variables commute:
\[
  x_i x_j = x_j x_i.
\]
These are the ordinary variables used for symmetric, quasisymmetric, and
polynomial bases throughout the site.

\textbf{Fermionic variables.}
Fermionic variables anticommute:
\[
  \theta_i\theta_j=-\theta_j\theta_i.
\]
In particular, over a field of characteristic zero, \(\theta_i^2=0\).
They appear in super and fermionic variants of symmetric-function and
\hyperref[diagonalHarmonics]{diagonal-harmonic} constructions.

\textbf{Plethystic notation.}
Plethystic notation is a compact way to substitute alphabets and virtual
alphabets into symmetric functions.  It is tied to the
\hyperref[plethysm]{plethysm} operation.


\subsection[dictionaryPartitionsTableaux]{Partitions, tableaux, and Catalan objects}

\textbf{Frobenius notation.}
The \hyperref[frobeniusCoordinates]{Frobenius notation} of a partition records
the arm and leg lengths from the diagonal boxes of its Young diagram.

\textbf{Plane partition.}
A \hyperref[planePartitions]{plane partition} is a two-dimensional array of
nonnegative integers with weakly decreasing rows and columns.

\textbf{Tamari lattice.}
The Tamari lattice is a partial order on Catalan objects such as binary trees
or Dyck paths.  It also appears in rational, cyclic, affine, and
\(\nu\)-Tamari variants; see the \hyperref[cspRationalCatalan]{rational
Catalan} and \hyperref[tamariInterval]{Tamari interval} discussions.


\subsection[dictionaryGeometryPolytopes]{Geometry and polytopes}

\textbf{Ehrhart polynomial.}
The \hyperref[ehrhartPolynomialDefinition]{Ehrhart polynomial} of a lattice
polytope \(P\) counts lattice points in dilates \(nP\).

\textbf{\(h^*\)-polynomial.}
The \hyperref[hStarPolynomial]{\(h^*\)-polynomial} is the numerator of the
Ehrhart series of a lattice polytope.  Its coefficients are always
nonnegative integers, even though the Ehrhart polynomial itself can have
negative coefficients.

\textbf{Hessenberg variety.}
A \hyperref[hessenbergVarietyDefinition]{Hessenberg variety} is a subvariety
of a \hyperref[flagVarietyDefinition]{flag variety} defined by a linear
operator and a \hyperref[hessenbergFunctionDefinition]{Hessenberg function}.
Regular semisimple Hessenberg varieties are connected to chromatic
quasisymmetric functions through Tymoczko's
\hyperref[tymoczkoDotAction]{dot action}.

\textbf{Schubert cell.}
A Schubert cell is one stratum in the standard cell decomposition of a flag
variety or Grassmannian.  In the Grassmannian, these cells are indexed by
partitions fitting inside a rectangle and are closely related to
\hyperref[grassmannianSchubertClasses]{Schubert classes}.
